\documentclass[11pt]{article}

\usepackage[a4paper,margin=1in]{geometry}
\usepackage{amsmath,amssymb}
\usepackage{hyperref}

\hypersetup{
    colorlinks=true,
    linkcolor=blue,
    citecolor=blue,
    urlcolor=blue
}

\title{Why the Minimal-Energy Velocity Is a Gradient}
\author{}
\date{}

\begin{document}

\maketitle

In Otto calculus, a velocity field $v$ generates a tangent vector at a density
$\rho$ through the continuity equation:
\[
\dot\rho=-\operatorname{div}(\rho v).
\]
Different velocity fields can generate the same $\dot\rho$. Indeed, if $w$
satisfies
\[
\operatorname{div}(\rho w)=0,
\]
then
\[
-\operatorname{div}(\rho(v+w))
=
-\operatorname{div}(\rho v).
\]
So $v$ and $v+w$ represent the same tangent vector.

\section*{Choosing the Minimal-Energy Representative}

The kinetic energy of a velocity field is
\[
\int_{\mathbb R^d}|v|^2\,d\rho.
\]
Among all velocity fields producing the same tangent vector $\dot\rho$, Otto
calculus chooses the one with minimal kinetic energy.

This is analogous to choosing the shortest vector among all representatives of
an equivalence class.

\section*{Why It Is a Gradient}

The gradient fields are orthogonal, in the weighted inner product
\[
\langle u,v\rangle_\rho
=
\int_{\mathbb R^d} u\cdot v\,d\rho,
\]
to the fields that do not move the density, namely fields $w$ satisfying
\[
\operatorname{div}(\rho w)=0.
\]
Indeed, for a smooth test potential $\phi$,
\[
\int_{\mathbb R^d} \nabla\phi\cdot w\,d\rho
=
\int_{\mathbb R^d} \nabla\phi\cdot \rho w\,dx
=
-\int_{\mathbb R^d}\phi\,\operatorname{div}(\rho w)\,dx
=0.
\]
Thus any velocity field can be decomposed formally as
\[
v=\nabla\phi+w,
\qquad
\operatorname{div}(\rho w)=0.
\]
The component $w$ does not affect the tangent vector, but it increases kinetic
energy:
\[
\int |v|^2\,d\rho
=
\int |\nabla\phi|^2\,d\rho
+
\int |w|^2\,d\rho.
\]
Therefore the minimal-energy representative is obtained by setting $w=0$:
\[
v=\nabla\phi.
\]

\section*{Interpretation}

Only the compressive part of the velocity field changes the density. The
weighted divergence-free part moves mass around without changing the density to
first order, so it is extra motion. Wasserstein geometry discards that extra
motion and keeps the gradient velocity.

\end{document}
